\section{Capítulo IV: Impacto socioeconómico y perspectivas futuras}
\subsection{Impacto de la RDNFO en la inclusión digital}
Antes de evaluar los efectos concretos de la Red Dorsal Nacional de Fibra Óptica (RDNFO) sobre los servicios públicos digitales, conviene precisar en qué medida esta infraestructura ha logrado traducirse en una reducción real de la brecha digital, más allá del despliegue físico del cableado. Al respecto, el Grupo de Trabajo Sectorial del Ministerio de Transportes y Comunicaciones (MTC) reconoció en su informe final que, pese a que la red se encuentra desplegada en la práctica totalidad de las capitales de provincia previstas, su nivel de aprovechamiento efectivo continúa siendo mínimo (MTC, 2022). Esto evidencia una disociación entre la meta de cobertura física, ya cumplida, y el objetivo sustantivo de la Ley Nro. 29904, que era masificar el acceso a la banda ancha.

En la misma línea, la tesis de maestría de Rosas Lozada (2021), elaborada en la Pontificia Universidad Católica del Perú (PUCP), constató que, a mayo de 2021, únicamente el 3.2\% de la capacidad instalada de la RDNFO y de los proyectos regionales de banda ancha se encontraba en uso. Aunque para abril de 2023 dicho aprovechamiento había subido levemente hasta un 7.8\%, este avance resulta marginal frente a la magnitud de la inversión realizada y confirma que el problema no reside en la ausencia de infraestructura de transporte, sino en la falta de mecanismos que conecten dicha infraestructura con la demanda real de los usuarios finales.

Este contraste se vuelve más notorio al comparar la escasa evolución del uso de la RDNFO con el dinamismo que ha mostrado el mercado de telecomunicaciones peruano en los últimos tres años. Según cifras del OSIPTEL basadas en mediciones de Ookla, la velocidad mediana de descarga de Internet fijo en el país pasó de 94.6 Mbps en diciembre de 2023 a 195.9 Mbps en diciembre de 2024, más del doble en solo doce meses, y continuó su ascenso hasta superar los 200 Mbps en abril de 2025, lo que ubicó al Perú en el segundo lugar de Sudamérica en esta variable, solo detrás de Chile (OSIPTEL, 2025). Sin embargo, conviene precisar que dicha mejora responde principalmente a la inversión de los operadores privados en sus propias redes de última milla y en tecnologías de mayor capacidad, y no al aprovechamiento de la infraestructura troncal estatal, cuyo nivel de uso permanece prácticamente estancado. En otras palabras, el país ha logrado avances notables en velocidad de conexión pese a la subutilización de la RDNFO, y no necesariamente gracias a ella.

Este hallazgo se refuerza al observar los datos más recientes del Instituto Nacional de Estadística e Informática (INEI), según los cuales el porcentaje de hogares con acceso a Internet a nivel nacional pasó de 58.4 \% a fines de 2024 a 60.1\% en el segundo trimestre de 2025 (INEI, 2025). No obstante, esta mejora agregada oculta una brecha persistente entre zonas urbanas y rurales; mientras que en Lima Metropolitana el acceso a Internet alcanzaba el 80.5\% de los hogares, en el área rural apenas llegaba al 23.6\% en el mismo periodo (INEI, 2025). Esta disparidad confirma, con cifras actualizadas, un diagnóstico que la propia CEPAL viene sosteniendo hasta sus informes más recientes: la infraestructura digital avanzada continúa siendo uno de los principales cuellos de botella de la región, pues aunque la conectividad móvil ha avanzado, persisten brechas profundas en el acceso equitativo a banda ancha fija de calidad, especialmente en zonas rurales (CEPAL, 2024). Aplicado al caso peruano, esto sugiere que la sola disponibilidad de infraestructura de transporte no basta para cerrar la brecha digital si no se articula con políticas específicas de estímulo a la demanda y de despliegue de última milla en las zonas de menor rentabilidad comercial, que es precisamente donde la RDNFO debería tener mayor incidencia.

\subsection{Impacto en los servicios públicos digitales}
\subsubsection{Teleeducación}
\subsubsection{Telesalud}
\subsubsection{Gobierno digital}
\subsection{Consecuencias socioeconómicas de la subutilización de la RDNFO}
\subsubsection{Consecuencias sociales}
\subsubsection{Consecuencias económicas}
\subsubsection{Consecuencias territoriales}
\subsection{Experiencias comparadas y lecciones para el Perú}
\subsection{Perspectivas futuras y propuestas de optimización}
\subsubsection{Flexibilización tarifaria}
\subsubsection{Alianzas público-privadas}
\subsubsection{Nuevos nodos y fortalecimiento de la red}
